% ============================================================
%  Notes on Tokenization, Perplexity, Cross-Entropy,
%  Scaling Laws, Transformers, and Emergent Reasoning
% ============================================================
\documentclass[11pt,a4paper]{article}

% ---- Layout ------------------------------------------------
\usepackage[top=1in, bottom=1in, left=1.1in, right=1.1in]{geometry}

% ---- Typography --------------------------------------------
\usepackage[T1]{fontenc}
\usepackage{lmodern}
\usepackage{microtype}

% ---- Math --------------------------------------------------
\usepackage{amsmath, amssymb, amsthm, mathtools, bm}

% ---- Colors & Graphics -------------------------------------
\usepackage[dvipsnames,table]{xcolor}
\usepackage{graphicx}
\usepackage{tikz}
\usepackage{xurl}

% ---- Tables ------------------------------------------------
\usepackage{booktabs}
\usepackage{array}
\usepackage{tabularx}

% ---- Lists -------------------------------------------------
\usepackage{enumitem}

% ---- Section Formatting ------------------------------------
\usepackage{titlesec}
\usepackage{titletoc}

% ---- Header / Footer ---------------------------------------
\usepackage{fancyhdr}

% ---- Boxes & Frames ----------------------------------------
\usepackage{tcolorbox}
\tcbuselibrary{skins, breakable, theorems}

% ---- Hyperlinks --------------------------------------------
\usepackage{hyperref}
\hypersetup{
  colorlinks = true,
  linkcolor  = RoyalBlue,
  citecolor  = RoyalBlue,
  urlcolor   = RoyalBlue,
  pdftitle   = {Notes on LLMs: Theory and Foundations},
  pdfauthor  = {},
}

% ---- Bibliography ------------------------------------------
\usepackage[numbers, sort&compress]{natbib}
\bibliographystyle{plainnat}

% ============================================================
%  COLOUR PALETTE
% ============================================================
\definecolor{primary}{RGB}{31, 73, 125}      % deep navy
\definecolor{accent}{RGB}{0, 112, 192}       % bright blue
\definecolor{boxbg}{RGB}{235, 243, 252}      % light blue tint
\definecolor{theorembg}{RGB}{240, 248, 240}  % light green tint
\definecolor{warningbg}{RGB}{255, 248, 230}  % light amber
\definecolor{rulecolor}{RGB}{0, 112, 192}
\definecolor{shaderow}{RGB}{240, 246, 255}

% ============================================================
%  SECTION FORMATTING
% ============================================================
\titleformat{\section}
  {\large\bfseries\color{primary}}
  {\thesection}{1em}{}
  [\vspace{-2pt}\textcolor{accent}{\rule{\linewidth}{0.6pt}}]

\titleformat{\subsection}
  {\normalsize\bfseries\color{accent}}
  {\thesubsection}{1em}{}

\titleformat{\subsubsection}
  {\normalsize\bfseries\itshape\color{primary!80}}
  {\thesubsubsection}{1em}{}

\titlespacing{\section}{0pt}{16pt}{6pt}
\titlespacing{\subsection}{0pt}{10pt}{4pt}

% ============================================================
%  HEADER / FOOTER
% ============================================================
\pagestyle{fancy}
\fancyhf{}
\setlength{\headheight}{22pt}
\fancyhead[L]{\small\color{primary}\textit{LLM Theory: Foundations}}
\fancyhead[R]{\small\color{primary}\thepage}
\fancyfoot[C]{\small\color{gray!60}\rule{0.4\linewidth}{0.4pt}}
\renewcommand{\headrulewidth}{0.4pt}
\renewcommand{\headrule}{\hbox to\headwidth{\color{accent}\leaders\hrule height \headrulewidth\hfill}}

% ============================================================
%  TCOLORBOX ENVIRONMENTS
% ============================================================

%% Key Result box (blue)
\newtcolorbox{keyresult}[1][]{
  enhanced, breakable,
  colback   = boxbg,
  colframe  = accent,
  arc       = 3pt,
  boxrule   = 0.8pt,
  left      = 8pt, right = 8pt, top = 6pt, bottom = 6pt,
  fonttitle = \bfseries\small\color{primary},
  title     = {#1},
}

%% Definition box (green tint)
\newtcolorbox{defbox}[1][]{
  enhanced, breakable,
  colback   = theorembg,
  colframe  = OliveGreen!70,
  arc       = 3pt,
  boxrule   = 0.8pt,
  left      = 8pt, right = 8pt, top = 6pt, bottom = 6pt,
  fonttitle = \bfseries\small\color{OliveGreen!80!black},
  title     = {#1},
}

%% Caution / remark box (amber)
\newtcolorbox{cautionbox}[1][]{
  enhanced, breakable,
  colback   = warningbg,
  colframe  = Goldenrod!80!black,
  arc       = 3pt,
  boxrule   = 0.8pt,
  left      = 8pt, right = 8pt, top = 6pt, bottom = 6pt,
  fonttitle = \bfseries\small\color{Goldenrod!70!black},
  title     = {#1},
}

%% Synthesis / summary box (grey tint)
\newtcolorbox{synthbox}[1][Summary]{
  enhanced, breakable,
  colback   = gray!7,
  colframe  = gray!50,
  arc       = 3pt,
  boxrule   = 0.8pt,
  left      = 8pt, right = 8pt, top = 6pt, bottom = 6pt,
  fonttitle = \bfseries\small\color{gray!60!black},
  title     = {#1},
}

% ============================================================
%  THEOREM ENVIRONMENTS
% ============================================================
\theoremstyle{definition}
\newtheorem{definition}{Definition}[section]
\newtheorem{remark}{Remark}[section]
\newtheorem{example}{Example}[section]

\theoremstyle{plain}
\newtheorem{proposition}{Proposition}[section]

% ============================================================
%  CUSTOM COMMANDS
% ============================================================
\newcommand{\R}{\mathbb{R}}
\newcommand{\E}{\mathbb{E}}
\newcommand{\KL}[2]{\mathrm{KL}\!\left(#1\,\|\,#2\right)}
\newcommand{\PPL}{\mathrm{PPL}}
\newcommand{\Hc}{H}
\newcommand{\tloss}{\mathcal{L}}

% ============================================================
%  TABLE OF CONTENTS FORMATTING
% ============================================================
\setcounter{tocdepth}{2}

% ============================================================
%  DOCUMENT
% ============================================================
\begin{document}

% ---- Title Page --------------------------------------------
\hypersetup{pageanchor=false}
\begin{titlepage}
  \pagecolor{primary}
  \color{white}
  \vspace*{4cm}
  \begin{center}
    {\Huge\bfseries Notes on Large Language Models}\\[0.6em]
    {\Large\itshape Theory and Foundations}\\[2.5em]
    \textcolor{accent!60!white}{\rule{0.5\linewidth}{1pt}}\\[1.8em]
    {\normalsize
      Tokenization\enspace$\cdot$\enspace
      Cross-Entropy \& Perplexity\enspace$\cdot$\enspace
      KL Divergence\\[0.4em]
      Scaling Laws\enspace$\cdot$\enspace
      Transformer Architecture\enspace$\cdot$\enspace
      Emergent Reasoning
    }\\[3em]
    {\small\color{gray!40} Self-study reference notes}
  \end{center}
  \vfill
  \begin{center}
    {\small\color{gray!40}\today}
  \end{center}
\end{titlepage}
\pagecolor{white}\color{black}

% ---- TOC ---------------------------------------------------
\pagenumbering{roman}
\tableofcontents
\newpage
\hypersetup{pageanchor=true}
\pagenumbering{arabic}

% ============================================================
\section{Introduction}
% ============================================================

These notes provide a mathematically grounded treatment of several core concepts
underlying modern large language models (LLMs).
The exposition is self-contained and proceeds from first principles,
covering the following topics in order:

\begin{enumerate}[leftmargin=2em, itemsep=2pt]
  \item \textbf{Tokenization} --- how raw text is mapped into discrete symbols,
        and why this choice affects both computational and statistical difficulty.
  \item \textbf{Language modeling objective} --- autoregressive factorization
        and negative log-likelihood training.
  \item \textbf{Cross-entropy and KL divergence} --- the information-theoretic
        interpretation of LLM training.
  \item \textbf{Perplexity} --- a practical measure of predictive uncertainty.
  \item \textbf{Scaling laws} --- empirical power-law relationships between
        scale and loss.
  \item \textbf{Transformer architecture} --- self-attention, feedforward layers,
        and expressive capacity.
  \item \textbf{Emergent reasoning} --- why structured reasoning can arise
        from next-token prediction at scale.
\end{enumerate}

\medskip
\noindent
\textbf{Notation.}\enspace
Throughout, let raw text be mapped into a sequence of discrete tokens
\[
  x \;\mapsto\; (t_1, t_2, \dots, t_n), \qquad t_i \in V,
\]
where $V$ is a finite vocabulary.
We write $t_{<i} := (t_1, \dots, t_{i-1})$ for the prefix preceding position $i$.
$P^*$ denotes the true (unknown) distribution of natural language text.
$P_\theta$ denotes the model's learned distribution, parameterised by $\theta$.

% ============================================================
\section{Tokenization}
% ============================================================

\subsection{Formal Definition}

\begin{defbox}[Definition: Tokenizer]
A \textit{tokenizer} is a mapping
\[
  T : \mathcal{X} \;\longrightarrow\; V^{*},
\]
from raw text space $\mathcal{X}$ to finite sequences over a discrete vocabulary $V$.
The output $T(x) = (t_1, \dots, t_n)$ is the \emph{token sequence} on which
language modeling is performed.
\end{defbox}

Language modeling therefore operates on $T(x)$, not on $x$ directly.
This indirection has far-reaching consequences.

\subsection{Computational Consequences}

For a Transformer, self-attention over a sequence of length $n$ costs approximately
\[
  O\!\left(n^2 d\right),
\]
where $d$ is the model's hidden dimension \cite{vaswani2017attention}.
A tokenizer that produces longer sequences therefore increases:
compute per forward pass, peak memory, inference latency, and
the effective horizon that must be handled in long-context settings.
Efficient tokenization directly reduces these costs.

\subsection{Statistical Consequences}

Tokenization changes the \emph{coordinate system} of the learning problem.
Specifically, it alters:
\begin{itemize}[leftmargin=2em, itemsep=2pt]
  \item the sequence length $n$,
  \item the vocabulary distribution $\Pr(t)$,
  \item the difficulty of each conditional $P^*(t_i \mid t_{<i})$,
  \item and the amount of data available to estimate each unit.
\end{itemize}

A rare token type is harder to model well than a frequent one,
so the tokenizer's choice of merge granularity affects statistical efficiency.

\subsection{Compression--Learnability Trade-off}

A good tokenizer must balance two competing objectives:

\begin{enumerate}[leftmargin=2em, itemsep=2pt]
  \item \textbf{Compression}: short token sequences reduce sequence length $n$.
  \item \textbf{Learnability}: token units must recur frequently enough for
        the model to estimate their conditionals reliably.
\end{enumerate}

This can be formalised through a \emph{minimum description length} (MDL) lens \cite{grunwald2007mdl}:
\[
  L(\text{text})
  \;\approx\;
  \underbrace{L(\text{token sequence})}_{\text{representation cost}}
  \;+\;
  \underbrace{L(\text{model needed to predict it})}_{\text{model cost}}.
\]
A tokenizer that minimises only the first term (very long tokens, small $n$)
may dramatically increase the second term because rare token types are hard to model.
The optimal tokenizer minimises the total description length.

\subsection{Byte Pair Encoding}

Byte Pair Encoding (BPE) \cite{sennrich2016bpe} is the most widely used subword
tokenization algorithm.
It begins with a character-level vocabulary and iteratively merges the most frequent
adjacent pair of symbols, until a target vocabulary size is reached.
Formally, the guiding heuristic is:
\[
  \text{frequent adjacent patterns} \;\longrightarrow\; \text{single merged token}.
\]
WordPiece \cite{schuster2012wordpiece} and Unigram LM \cite{kudo2018sentencepiece}
are alternative approaches that optimise slightly different objectives.

\begin{keyresult}[Key insight: tokenization as compression]
BPE achieves compression by assigning short codes (single tokens) to frequent patterns,
analogous to Huffman coding in information theory.
The expected code length under BPE is close to the empirical entropy of the
byte-pair distribution on the training corpus.
\end{keyresult}

\subsection{Tokenization Biases}

Tokenization can introduce systematic biases that affect fairness and performance:

\begin{table}[ht]
\centering
\small
\renewcommand{\arraystretch}{1.35}
\rowcolors{2}{shaderow}{white}
\begin{tabularx}{\linewidth}{@{} l X @{}}
  \toprule
  \rowcolor{primary!15}
  \textbf{Bias type} & \textbf{Description} \\
  \midrule
  Frequency bias      & Frequent words receive shorter encodings than rare words, making rare words more expensive to model. \\
  Language bias       & Languages underrepresented in tokenizer training may be fragmented into more tokens per word, increasing cost and degrading performance \cite{rust2021good}. \\
  Domain bias         & Code, mathematics, biological sequences, and financial identifiers may be split inefficiently across tokens. \\
  Morphological bias  & Morphologically rich languages (Finnish, Turkish, Arabic) produce much longer sequences, compounding all the above. \\
  Numeric bias        & Numbers may be split character-by-character, making arithmetic harder \cite{nogueira2021investigating}. \\
  \bottomrule
\end{tabularx}
\caption{Common biases induced by subword tokenization.}
\label{tab:tokenization-biases}
\end{table}

% ============================================================
\section{The Language Modeling Objective}
% ============================================================

\subsection{Autoregressive Factorization}

By the chain rule of probability, any joint distribution over a sequence
$(t_1, \dots, t_n)$ factorises as:
\[
  P(t_1, \dots, t_n)
  = \prod_{i=1}^{n} P(t_i \mid t_{<i}).
\]

An \textit{autoregressive language model} parameterises each factor:
\[
  P_\theta(t_1, \dots, t_n)
  = \prod_{i=1}^{n} P_\theta(t_i \mid t_{<i}).
\]

This factorisation is exact (no independence assumption), causal (each prediction
conditions only on the past), and tractable (each factor is a distribution over
the finite vocabulary $V$).

\subsection{Training Objective}

Maximum likelihood estimation (MLE) over a training corpus $\mathcal{D}$
is equivalent to minimising the \emph{negative log-likelihood} (NLL):
\[
  \tloss(\theta)
  = -\sum_{i=1}^{n} \log P_\theta(t_i \mid t_{<i}).
\]

Averaged over the empirical data distribution $\hat{P}$, this becomes
\[
  \tloss(\theta)
  = -\,\E_{(t_1,\dots,t_n)\sim \hat{P}}
    \!\left[\sum_{i=1}^{n} \log P_\theta(t_i \mid t_{<i})\right].
\]

\begin{remark}
  MLE training with an autoregressive factorisation is also called
  \emph{teacher forcing}: the model is always conditioned on ground-truth
  prefixes during training, not on its own previous predictions.
  This creates a \emph{train-test discrepancy} (exposure bias) that can
  affect generation quality at test time \cite{bengio2015scheduled}.
\end{remark}

% ============================================================
\section{Cross-Entropy and KL Divergence}
% ============================================================

\subsection{Cross-Entropy}

Let $P^*(t_i \mid t_{<i})$ denote the true conditional distribution and
$P_\theta(t_i \mid t_{<i})$ the model's approximation.
The \emph{cross-entropy} between $P^*$ and $P_\theta$ is defined as:
\[
  \Hc(P^*, P_\theta)
  = -\,\E_{P^*}\!\left[\log P_\theta\right].
\]

The NLL training loss is precisely the empirical estimate of this cross-entropy.

\subsection{The KL Decomposition}

A fundamental identity in information theory is:

\begin{keyresult}[Cross-entropy decomposition]
\[
  \Hc(P^*, P_\theta)
  \;=\;
  \underbrace{\Hc(P^*)}_{\text{irreducible entropy}}
  \;+\;
  \underbrace{\KL{P^*}{P_\theta}}_{\text{model excess}},
\]
where $\Hc(P^*) = -\E_{P^*}[\log P^*]$ is the entropy of the true distribution
and $\KL{P^*}{P_\theta} = \E_{P^*}\!\left[\log \frac{P^*}{P_\theta}\right] \geq 0$.
\end{keyresult}

\noindent
Since $\Hc(P^*)$ does not depend on $\theta$, minimising cross-entropy is
\emph{exactly equivalent} to minimising the KL divergence from $P^*$ to $P_\theta$:
\[
  \min_\theta\;\Hc(P^*, P_\theta)
  \;\Longleftrightarrow\;
  \min_\theta\;\KL{P^*}{P_\theta}.
\]

\subsection{Interpretation}

\begin{keyresult}[Information-theoretic view of LLM training]
Training an LLM is equivalent to finding the closest approximation $P_\theta$
to the true distribution of language $P^*$, in the sense of KL divergence.
The irreducible entropy $\Hc(P^*)$ represents the intrinsic unpredictability
of natural language and sets a fundamental floor on achievable loss.
\end{keyresult}

\noindent
The quantity $\KL{P^*}{P_\theta}$ measures the information lost per token
when using the model $P_\theta$ instead of the true distribution $P^*$.
It equals zero if and only if $P_\theta = P^*$ almost everywhere.

\subsection{Bits-Per-Character and Other Units}

Cross-entropy is often reported in \emph{nats} (natural log) or \emph{bits}
(log base 2). The conversion is:
\[
  \text{bits-per-token}
  = \frac{\Hc(P^*, P_\theta)}{\ln 2}.
\]
Bits-per-character (BPC) is also common, normalising by the average number of
characters per token.

% ============================================================
\section{Perplexity}
% ============================================================

\subsection{Definition}

\begin{defbox}[Definition: Perplexity]
The \emph{perplexity} of a model $P_\theta$ on a token sequence $(t_1, \dots, t_n)$
is defined as:
\[
  \PPL
  = \exp\!\left(-\frac{1}{n}\sum_{i=1}^{n} \log P_\theta(t_i \mid t_{<i})\right)
  = e^{\,\ell},
\]
where $\ell$ is the mean NLL per token.
\end{defbox}

\subsection{Interpretation}

Perplexity admits several complementary interpretations:

\begin{enumerate}[leftmargin=2em, itemsep=2pt]
  \item \textbf{Geometric mean inverse probability.}
        $\PPL = \left(\prod_{i=1}^n P_\theta(t_i \mid t_{<i})\right)^{-1/n}$.
        This is the geometric mean of the reciprocal probabilities assigned
        to the correct next token.
  \item \textbf{Effective branching factor.}
        Heuristically, $\PPL \approx k$ means the model behaves as if it is
        choosing uniformly among $k$ equally likely continuations at each step.
  \item \textbf{Exponential of cross-entropy.}
        Since the mean NLL estimates $\Hc(P^*, P_\theta)$, we have
        $\PPL \approx e^{\Hc(P^*, P_\theta)}$.
\end{enumerate}

Lower perplexity indicates better predictive performance.
A model that perfectly predicted every token would have $\PPL = 1$.
A model that assigned uniform probability over $|V|$ tokens at every step would
have $\PPL = |V|$ (often $\sim 50{,}000$).

\subsection{Dependence on Tokenization}

\begin{cautionbox}[Caution: cross-tokenizer comparisons]
Perplexity values are \textbf{not directly comparable} across different tokenizers.
A tokenizer that splits text into more tokens (larger $n$) will generally yield
lower per-token perplexity for the same model, simply because each prediction
task is easier.
Normalisation by character count (bits-per-character) is necessary for
fair cross-tokenizer comparison.
\end{cautionbox}

% ============================================================
\section{How Tokenization Affects Loss}
% ============================================================

Consider a concept represented as a single token $u$ under tokenizer $T_1$,
but split into $m$ subtokens $(s_1, \dots, s_m)$ under tokenizer $T_2$.

The loss contributions are:
\[
  \tloss_{T_1} = -\log P_\theta(u \mid \text{context}),
  \qquad
  \tloss_{T_2} = -\sum_{j=1}^{m} \log P_\theta(s_j \mid s_{<j}, \text{context}).
\]

Under $T_2$, the model must predict $m$ tokens instead of one.
While each individual prediction may be easier (smaller effective vocabulary
at each step), the model must also maintain coherence across $m$ positions
and attend over a larger context window.

This affects:
\begin{itemize}[leftmargin=2em, itemsep=2pt]
  \item \textbf{Optimisation difficulty}: more gradient steps touch a single concept.
  \item \textbf{Attention burden}: dependencies span more positions.
  \item \textbf{Statistical efficiency}: each subtoken individually has less
        informational content, requiring more data to estimate its distribution well.
  \item \textbf{Final model quality}: tasks requiring exact reconstruction
        of multi-token units (e.g., arithmetic, identifiers) become harder.
\end{itemize}

% ============================================================
\section{Scaling Laws}
% ============================================================

\subsection{Empirical Form}

A central empirical finding \cite{kaplan2020scaling, hoffmann2022training} is
that language-model loss follows approximate power laws in scale:

\begin{keyresult}[Scaling law (Kaplan et al.)]
\[
  \tloss(N) \;\approx\; A\,N^{-\alpha} + B,
\]
where $N$ may denote model parameters, training tokens, or compute (FLOPs);
$A > 0$ and $\alpha > 0$ are empirically fitted constants;
and $B > 0$ is the \emph{irreducible loss floor} corresponding to $\Hc(P^*)$.
\end{keyresult}

\noindent
Typical values from the literature are $\alpha \approx 0.07$--$0.1$ for model
size and $\alpha \approx 0.27$--$0.35$ for data \cite{kaplan2020scaling}.

\subsection{The Chinchilla Scaling Laws}

Hoffmann et al.\ \cite{hoffmann2022training} (``Chinchilla'') established a
critical revision of earlier scaling prescriptions.
Given a fixed compute budget $C$ (measured in FLOPs), the optimal allocation is:
\[
  N^* \;\propto\; C^{0.5}, \qquad D^* \;\propto\; C^{0.5},
\]
implying that the number of training tokens $D$ should scale roughly \emph{linearly}
with the number of parameters $N$.
This led to the practical rule of thumb:
\[
  D^* \approx 20 \times N.
\]

Many models prior to Chinchilla were significantly \emph{undertrained} relative
to their parameter count.

\subsection{Why Scaling Helps}

Natural language contains structure at many levels simultaneously:
\begin{itemize}[leftmargin=2em, itemsep=2pt]
  \item local syntax and morphology,
  \item long-range semantic coherence,
  \item world knowledge and factual regularities,
  \item discourse and pragmatic patterns,
  \item latent reasoning chains.
\end{itemize}

Larger models have sufficient \emph{capacity} to represent more of these patterns.
Larger datasets \emph{expose} the model to more of them.
More compute allows optimisation to \emph{exploit} this capacity.
Together, these push $P_\theta \to P^*$.

\subsection{Emergent Capabilities}

Beyond smooth reductions in loss, scaling is associated with \emph{emergent capabilities}:
qualitative abilities that appear to arise discontinuously above certain thresholds
\cite{wei2022emergent}.
Examples include few-shot in-context learning, multi-step arithmetic, and
chain-of-thought reasoning.
The emergence phenomenon is debated: some analyses suggest it is an artefact of
non-linear evaluation metrics rather than a true phase transition
\cite{schaeffer2023emergent}.

% ============================================================
\section{Transformer Architecture}
% ============================================================

\subsection{Input Representation}

Each token $t_i$ is first mapped to a $d$-dimensional embedding vector $\mathbf{e}_i \in \R^d$,
and a \emph{positional encoding} $\mathbf{p}_i \in \R^d$ is added:
\[
  \mathbf{x}_i^{(0)} = \mathbf{e}_i + \mathbf{p}_i.
\]

Original Transformers used sinusoidal positional encodings \cite{vaswani2017attention};
modern LLMs typically use Rotary Position Embeddings (RoPE) \cite{su2024rope}
or ALiBi \cite{press2022alibi}, which generalise better to sequence lengths
beyond those seen during training.

\subsection{Self-Attention}

At each layer $\ell$, the input matrix $X \in \R^{n \times d}$ is linearly
projected into queries, keys, and values:
\[
  Q = X W_Q, \quad K = X W_K, \quad V = X W_V,
  \qquad W_Q, W_K, W_V \in \R^{d \times d_k}.
\]

The scaled dot-product attention output is:
\[
  \mathrm{Attention}(Q, K, V)
  = \mathrm{softmax}\!\left(\frac{QK^\top}{\sqrt{d_k}}\right)V.
\]

The division by $\sqrt{d_k}$ prevents the dot products from growing large in
magnitude when $d_k$ is large, which would saturate the softmax and cause
vanishing gradients.

\subsubsection{Causal Masking}

For autoregressive generation, a \emph{causal mask} is applied to the attention
matrix, setting all entries at positions $(i, j)$ with $j > i$ to $-\infty$
before the softmax.
This ensures that $P_\theta(t_i \mid t_{<i})$ cannot attend to future tokens:
\[
  A_{ij} = \begin{cases}
    \dfrac{\mathbf{q}_i \cdot \mathbf{k}_j}{\sqrt{d_k}} & j \leq i, \\[6pt]
    -\infty & j > i.
  \end{cases}
\]

\subsubsection{Multi-Head Attention}

In practice, attention is computed in $H$ parallel \emph{heads}, each with
its own projection matrices $W_Q^{(h)}, W_K^{(h)}, W_V^{(h)} \in \R^{d \times d_k}$
where $d_k = d/H$:
\[
  \mathrm{MultiHead}(Q, K, V)
  = \mathrm{Concat}\!\left(\mathrm{head}_1, \dots, \mathrm{head}_H\right) W_O,
\]
\[
  \mathrm{head}_h = \mathrm{Attention}\!\left(QW_Q^{(h)}, KW_K^{(h)}, VW_V^{(h)}\right).
\]

Multiple heads allow the model to jointly attend to information from different
representation subspaces at different positions.

\subsection{Position-wise Feedforward Layer}

Each Transformer layer also contains a two-layer feedforward network applied
independently to each position:
\[
  \mathrm{FFN}(\mathbf{x}) = \sigma\!\left(\mathbf{x} W_1 + \mathbf{b}_1\right) W_2 + \mathbf{b}_2,
\]
where $\sigma$ is typically GELU \cite{hendrycks2016gelu} or SwiGLU \cite{shazeer2020glu}
in modern LLMs, and the inner dimension is expanded (commonly $4d$).

\subsection{Layer Normalization and Residual Connections}

Each sub-layer (attention or FFN) is wrapped with a \emph{residual connection}
and \emph{layer normalization} \cite{ba2016layer}:
\[
  \mathbf{x}^{(\ell+1)} = \mathrm{LayerNorm}\!\left(\mathbf{x}^{(\ell)} + \mathrm{SubLayer}\!\left(\mathbf{x}^{(\ell)}\right)\right).
\]

Modern LLMs commonly use \emph{pre-normalization} (applying LayerNorm before
the sub-layer) for improved training stability:
\[
  \mathbf{x}^{(\ell+1)} = \mathbf{x}^{(\ell)} + \mathrm{SubLayer}\!\left(\mathrm{LayerNorm}\!\left(\mathbf{x}^{(\ell)}\right)\right).
\]

\subsection{Why Transformers Are Expressive}

Transformers achieve high expressiveness through the interplay of:
\begin{enumerate}[leftmargin=2em, itemsep=2pt]
  \item \textbf{Dynamic, input-dependent routing}: attention weights
        $\mathrm{softmax}(QK^\top/\sqrt{d_k})$ depend on the actual content,
        not fixed connections.
  \item \textbf{Nonlinear transformation}: the FFN layers introduce pointwise
        nonlinearity at each position.
  \item \textbf{Depth and composition}: stacking $L$ layers allows successive
        refinement of representations.
  \item \textbf{Residual learning}: skip connections preserve gradient flow and
        allow layers to learn incremental refinements.
\end{enumerate}

It has been shown that Transformers are Turing-complete under certain conditions
\cite{perez2021turing}, and universal approximators of sequence-to-sequence functions
of bounded length \cite{yun2020transformers}.

\begin{keyresult}[Representational capacity of Transformers]
A sufficiently deep and wide Transformer can approximate any measurable
sequence-to-sequence function on sequences of bounded length,
given sufficient precision in its weights.
\end{keyresult}

% ============================================================
\section{What an LLM Is Learning}
% ============================================================

\subsection{Surface and Deep Views}

At the surface level, an LLM learns the next-token conditional:
\[
  P_\theta(t_i \mid t_{<i}).
\]

However, accurately modelling this distribution requires the model to implicitly
learn a vast amount of structure. To predict the next word in
\emph{``The capital of France is \underline{\hspace{1em}}''},
the model must encode the factual association between France and Paris.
To predict the next line in a proof, it must capture the relevant logical step.

\subsection{Latent Representation Learning}

Internally, the model builds a hierarchy of \emph{hidden representations}:
\[
  \mathbf{h}_i^{(\ell)} \in \R^d, \quad i = 1, \dots, n, \quad \ell = 1, \dots, L.
\]

Empirically, these encode increasingly abstract information with depth:
early layers capture token-level and syntactic features, middle layers
semantic and world-knowledge features, and late layers task-specific and
discourse-level features \cite{tenney2019bert, elhage2021mathematical}.

\subsection{The Bitter Lesson and Scale}

\citet{sutton2019bitter} argues that approaches leveraging computation
tend to outperform approaches that hand-craft structure.
LLMs embody this: rather than encoding specific linguistic knowledge,
they learn it from data at scale, and the learned representations often
surpass hand-engineered alternatives.

% ============================================================
\section{Why Reasoning Emerges from Next-Token Prediction}
% ============================================================

\subsection{Core Mechanism}

\begin{keyresult}[Emergence of reasoning]
Many next-token predictions are only possible if the model internalises
latent computational structure. Training pressure therefore incentivises the
development of internal reasoning capabilities---not because reasoning is
directly supervised, but because it is \emph{necessary for accurate prediction}.
\end{keyresult}

\subsection{Formal Perspective}

The training objective is:
\[
  \min_\theta -\,\E_{P^*}\!\left[\log P_\theta(\text{text})\right].
\]

The training corpus contains extensive reasoning traces:
mathematical proofs, worked examples, explanations, step-by-step tutorials,
causal arguments, and logical discourse.
Minimising next-token loss incentivises the model to reproduce these patterns,
and reproducing them faithfully requires the model to simulate the underlying
reasoning process.

\subsection{Marginalisation Over Latent Computation}

Let $z$ denote a latent computation (e.g., an intermediate reasoning step)
that is not explicitly present in the token sequence.
The model's prediction can be written as:
\[
  P_\theta(t_{i+1} \mid c)
  = \sum_{z} P_\theta(t_{i+1} \mid z, c)\,P_\theta(z \mid c).
\]

Even though $z$ is never directly supervised, the model may develop internal
states that approximate this marginalisation, because doing so reduces loss.

\subsection{In-Context Learning}

A striking emergent capability is \emph{in-context learning} (ICL)
\cite{brown2020gpt3}: given a few examples in the prompt, the model adapts
its behaviour without any weight update.

A Bayesian interpretation \cite{xie2021explanation} suggests that ICL implements
implicit posterior inference over the task.
Because the training data contains tasks with different structures, the model
learns a prior over tasks, and uses the in-context examples as evidence to
update this prior at inference time.

\subsection{Chain-of-Thought Reasoning}

Chain-of-thought (CoT) prompting \cite{wei2022cot} significantly improves
performance on multi-step reasoning tasks by prompting the model to generate
explicit intermediate steps before the final answer.

Instead of computing
\[
  P_\theta(a \mid c),
\]
the model instead generates a reasoning trajectory:
\[
  P_\theta(a \mid c)
  \;\approx\;
  \sum_{r_1, \dots, r_m}
  P_\theta(a \mid r_1, \dots, r_m, c)\,
  \prod_{j=1}^{m} P_\theta(r_j \mid r_{<j}, c).
\]

Making latent steps explicit provides the model with more positions in which
to carry out intermediate computation, effectively increasing the depth of
computation available at inference time.

\begin{remark}
  \citet{feng2023towards} showed that CoT can be necessary: there exist tasks
  for which Transformers without scratchpad intermediate tokens cannot represent
  the correct function, but Transformers \emph{with} CoT can.
\end{remark}

\subsection{Limits: Perplexity Does Not Imply Reasoning}

\begin{cautionbox}[Important caveat]
Low perplexity does \textbf{not} automatically imply strong reasoning ability.
Perplexity measures \emph{average} next-token prediction accuracy, but
reasoning requires:
\begin{itemize}[leftmargin=1.5em, itemsep=1pt]
  \item multi-step global consistency,
  \item stable intermediate state tracking,
  \item robust generalisation beyond training distributions,
  \item and often exact correctness rather than high-probability prediction.
\end{itemize}
Therefore: $\;\text{low perplexity} \;\not\Rightarrow\; \text{reliable reasoning}$.

Nevertheless, lower perplexity tends to correlate with better reasoning at scale,
because both arise from a more accurate approximation of $P^*$.
\end{cautionbox}

% ============================================================
\section{Overall Synthesis}
% ============================================================

The following table summarises the relationships between the core concepts
developed in these notes.

\begin{table}[ht]
\centering
\small
\renewcommand{\arraystretch}{1.4}
\rowcolors{2}{shaderow}{white}
\begin{tabularx}{\linewidth}{@{} l X @{}}
  \toprule
  \rowcolor{primary!15}
  \textbf{Concept} & \textbf{Role and key formula} \\
  \midrule
  Tokenization
    & Maps text to discrete tokens: $T:\mathcal{X}\to V^*$.
      Affects sequence length, vocabulary, and modelling difficulty. \\
  Autoregressive factorisation
    & $P_\theta(t_1,\dots,t_n) = \prod_i P_\theta(t_i\mid t_{<i})$.
      Exact, causal, tractable. \\
  Training loss
    & Minimise NLL: $\tloss(\theta) = -\sum_i \log P_\theta(t_i\mid t_{<i})$. \\
  Cross-entropy
    & $\Hc(P^*,P_\theta) = \Hc(P^*) + \KL{P^*}{P_\theta}$.
      Training minimises the KL term. \\
  Perplexity
    & $\PPL = e^{\,\ell}$.
      Effective branching factor; not comparable across tokenizers. \\
  Scaling laws
    & $\tloss(N)\approx AN^{-\alpha}+B$.
      Power-law improvement; Chinchilla-optimal: $D^*\approx 20N$. \\
  Transformer
    & $\mathrm{Attention}(Q,K,V)=\mathrm{softmax}\!\left(\frac{QK^\top}{\sqrt{d_k}}\right)V$.
      Dynamic routing + nonlinear composition. \\
  Emergent reasoning
    & Forced by prediction accuracy on structured text; enhanced by chain-of-thought
      and in-context learning. \\
  \bottomrule
\end{tabularx}
\caption{Summary of core concepts.}
\label{tab:synthesis}
\end{table}

\begin{synthbox}[Grand synthesis]
Modern LLMs are parameterised approximations of $P^*$, the true distribution of
human-generated text.
Tokenization defines the discrete representation space; autoregressive training
minimises KL divergence to $P^*$; perplexity quantifies remaining uncertainty;
and scaling laws describe how this uncertainty decreases with compute.
Transformers provide the expressive machinery to represent rich, context-dependent
functions; and reasoning emerges because accurate next-token prediction on
large-scale, structured human text necessitates the internalisation of latent
computational and inferential patterns.
\end{synthbox}

% ============================================================
\section{Conclusion}
% ============================================================

These notes have developed a unified, mathematically precise account of the
core ideas underlying modern LLMs.
The central thread is the information-theoretic framing: training is KL minimisation,
perplexity is exponential cross-entropy, and scaling moves $P_\theta$ closer to $P^*$.
Transformers are the architectural vehicle for this approximation, and their
expressiveness---combined with the implicit supervision provided by structured
human text---is sufficient to give rise to emergent capabilities including
in-context learning and chain-of-thought reasoning.

Several important directions are not covered here, including:
alignment and RLHF \cite{ouyang2022instructgpt},
efficient inference (speculative decoding, quantisation),
mixture-of-experts architectures \cite{fedus2022switch},
and mechanistic interpretability \cite{elhage2021mathematical}.
These are natural next topics for further study.

% ============================================================
%  REFERENCES
% ============================================================
\newpage
\begin{thebibliography}{99}

\bibitem{vaswani2017attention}
  Vaswani, A., Shazeer, N., Parmar, N., Uszkoreit, J., Jones, L., Gomez, A.\ N.,
  Kaiser, \L., and Polosukhin, I.
  \newblock Attention is all you need.
  \newblock \textit{Advances in Neural Information Processing Systems (NeurIPS)}, 30, 2017.

\bibitem{kaplan2020scaling}
  Kaplan, J., McCandlish, S., Henighan, T., Brown, T.\ B., Chess, B., Child, R.,
  Gray, S., Radford, A., Wu, J., and Amodei, D.
  \newblock Scaling laws for neural language models.
  \newblock \textit{arXiv:2001.08361}, 2020.

\bibitem{hoffmann2022training}
  Hoffmann, J., Borgeaud, S., Mensch, A., Buchatskaya, E., Cai, T., Rutherford, E.,
  \textit{et al.}
  \newblock Training compute-optimal large language models.
  \newblock \textit{Advances in Neural Information Processing Systems (NeurIPS)}, 35, 2022.

\bibitem{brown2020gpt3}
  Brown, T.\ B., Mann, B., Ryder, N., Subbiah, M., Kaplan, J., Dhariwal, P.,
  \textit{et al.}
  \newblock Language models are few-shot learners.
  \newblock \textit{Advances in Neural Information Processing Systems (NeurIPS)}, 33, 2020.

\bibitem{wei2022emergent}
  Wei, J., Tay, Y., Bommasani, R., Raffel, C., Zoph, B., Borgeaud, S.,
  \textit{et al.}
  \newblock Emergent abilities of large language models.
  \newblock \textit{Transactions on Machine Learning Research}, 2022.

\bibitem{schaeffer2023emergent}
  Schaeffer, R., Miranda, B., and Koyejo, S.
  \newblock Are emergent abilities of large language models a mirage?
  \newblock \textit{Advances in Neural Information Processing Systems (NeurIPS)}, 36, 2023.

\bibitem{wei2022cot}
  Wei, J., Wang, X., Schuurmans, D., Bosma, M., Ichter, B., Xia, F.,
  \textit{et al.}
  \newblock Chain-of-thought prompting elicits reasoning in large language models.
  \newblock \textit{Advances in Neural Information Processing Systems (NeurIPS)}, 35, 2022.

\bibitem{sennrich2016bpe}
  Sennrich, R., Haddow, B., and Birch, A.
  \newblock Neural machine translation of rare words with subword units.
  \newblock \textit{Proceedings of ACL}, 2016.

\bibitem{kudo2018sentencepiece}
  Kudo, T. and Richardson, J.
  \newblock SentencePiece: A simple and language independent subword tokenizer
  and detokenizer for neural text processing.
  \newblock \textit{Proceedings of EMNLP (System Demonstrations)}, 2018.

\bibitem{schuster2012wordpiece}
  Schuster, M. and Nakamura, K.
  \newblock Japanese and Korean voice search.
  \newblock \textit{IEEE International Conference on Acoustics, Speech and
  Signal Processing (ICASSP)}, 2012.

\bibitem{rust2021good}
  Rust, P., Pfeiffer, J., Vulic, I., Ruder, S., and Gurevych, I.
  \newblock How good is your tokenizer? On the monolingual performance of
  multilingual language models.
  \newblock \textit{Proceedings of ACL}, 2021.

\bibitem{nogueira2021investigating}
  Nogueira, R., Jiang, Z., and Lin, J.
  \newblock Investigating the limitations of Transformers with simple arithmetic tasks.
  \newblock \textit{arXiv:2102.13019}, 2021.

\bibitem{bengio2015scheduled}
  Bengio, S., Vinyals, O., Jaitly, N., and Shazeer, N.
  \newblock Scheduled sampling for sequence prediction with recurrent neural networks.
  \newblock \textit{Advances in Neural Information Processing Systems (NeurIPS)}, 28, 2015.

\bibitem{grunwald2007mdl}
  Gr\"{u}nwald, P.
  \newblock \textit{The Minimum Description Length Principle}.
  \newblock MIT Press, 2007.

\bibitem{su2024rope}
  Su, J., Ahmed, M., Lu, Y., Pan, S., Bo, W., and Liu, Y.
  \newblock RoFormer: Enhanced Transformer with Rotary Position Embedding.
  \newblock \textit{Neurocomputing}, 568, 2024.

\bibitem{press2022alibi}
  Press, O., Smith, N.\ A., and Lewis, M.
  \newblock Train short, test long: Attention with linear biases enables
  input length extrapolation.
  \newblock \textit{Proceedings of ICLR}, 2022.

\bibitem{hendrycks2016gelu}
  Hendrycks, D. and Gimpel, K.
  \newblock Gaussian error linear units (GELUs).
  \newblock \textit{arXiv:1606.08415}, 2016.

\bibitem{shazeer2020glu}
  Shazeer, N.
  \newblock GLU variants improve Transformer.
  \newblock \textit{arXiv:2002.05202}, 2020.

\bibitem{ba2016layer}
  Ba, J.\ L., Kiros, J.\ R., and Hinton, G.\ E.
  \newblock Layer normalization.
  \newblock \textit{arXiv:1607.06450}, 2016.

\bibitem{perez2021turing}
  P\'{e}rez, J., Barcel\'{o}, P., and Marinkovic, J.
  \newblock Attention is Turing-complete.
  \newblock \textit{Journal of Machine Learning Research}, 22(75):1--35, 2021.

\bibitem{yun2020transformers}
  Yun, C., Bhojanapalli, S., Rawat, A.\ S., Reddi, S.\ J., and Kumar, S.
  \newblock Are Transformers universal approximators of sequence-to-sequence functions?
  \newblock \textit{Proceedings of ICLR}, 2020.

\bibitem{tenney2019bert}
  Tenney, I., Das, D., and Pavlick, E.
  \newblock BERT rediscovers the classical NLP pipeline.
  \newblock \textit{Proceedings of ACL}, 2019.

\bibitem{elhage2021mathematical}
  Elhage, N., Nanda, N., Olsson, C., Henighan, T., Joseph, N., Mann, B.,
  \textit{et al.}
  \newblock A mathematical framework for Transformer circuits.
  \newblock \textit{Transformer Circuits Thread}, 2021.

\bibitem{xie2021explanation}
  Xie, S.\ M., Raghunathan, A., Liang, P., and Ma, T.
  \newblock An explanation of in-context learning as implicit Bayesian inference.
  \newblock \textit{Proceedings of ICLR}, 2022.

\bibitem[Sutton(2019)]{sutton2019bitter}
  Sutton, R.
  \newblock The bitter lesson.
  \newblock \textit{Incomplete Ideas (blog)}, 2019.
  \newblock \url{http://www.incompleteideas.net/IncIdeas/BitterLesson.html}.

\bibitem[Feng et~al.(2023)Feng, Zhang, Gu, Ye, He, and Wang]{feng2023towards}
  Feng, G., Zhang, B., Gu, Y., Ye, H., He, D., and Wang, L.
  \newblock Towards revealing the mystery behind chain of thought: A theoretical perspective.
  \newblock \textit{Advances in Neural Information Processing Systems (NeurIPS)}, 36, 2023.

\bibitem{ouyang2022instructgpt}
  Ouyang, L., Wu, J., Jiang, X., Almeida, D., Wainwright, C., Mishkin, P.,
  \textit{et al.}
  \newblock Training language models to follow instructions with human feedback.
  \newblock \textit{Advances in Neural Information Processing Systems (NeurIPS)}, 35, 2022.

\bibitem{fedus2022switch}
  Fedus, W., Zoph, B., and Shazeer, N.
  \newblock Switch Transformers: Scaling to trillion parameter models with simple
  and efficient sparsity.
  \newblock \textit{Journal of Machine Learning Research}, 23(120):1--39, 2022.

\end{thebibliography}

\end{document}
